\chapter{Conclusions}
\label{chap:conclusions}

This project set out to build a low-cost edge-processing telemetry unit for the TUKZIE Rev 0 electric cargo trike and to present its results safely to the driver. This chapter answers each research question from the evidence in Chapter~\ref{chap:results}, distinguishing what has been shown from what still depends on the vehicle.

\paragraph{RQ1: accuracy and repeatability of the dynamic-response measurement.} The measurement chain is shown to be sound on the bench: two inertial sensors are sampled at a fixed 200\,Hz with no lost samples, including with every other subsystem running; each is calibrated against gravity, which corrects a measured 8\% disagreement between nominally identical units; and the ride features and the dual-sensor fusion are computed correctly. How accurately and repeatably this chain characterises the trike's response to the road and the powertrain has not been measured, because no data have yet been recorded on the moving vehicle. RQ1 is therefore answered for the instrument but not for the vehicle.

\paragraph{RQ2: presenting safety-critical information on the existing dashboard.} A sparse presentation was implemented on the vehicle's own dashboard rather than a second, head-up display: a camera page with the live view and active alerts, a Sensors page with distances and fused hazards, and a banner over every page for an immediate hazard, all showing only live data and marking missing data as absent. On the bench the alerts reach the page within half a second of opening and the touch panel adds no measurable delay. Whether a driver can read this within the 2\,s glance guidance has not been measured, so RQ2 is answered for the information hierarchy and update strategy but not for the driver's glance time.

\paragraph{RQ3: latency, data loss and recovery.} On the bench, acquisition loses no samples; the only recurring pause, from writing the log to flash, is bounded to one 33 to 40\,ms gap a minute; the camera produces a result a median of 102\,ms after the sensor is read; every publish succeeded; the unit reconnects about 5\,s after coverage returns; and the local log is recovered completely. RQ3 is answered for the bench, with the same measures still to be repeated on the moving vehicle.

\paragraph{RQ4: features that separate road from powertrain vibration.} The time-domain features are implemented and verified against a synthetic signal, and the method for separating the two sources, using the motor's Hall pulses as a speed reference, is designed. Without vehicle data and the wired Hall tap, RQ4 is not answered.

\paragraph{Overall.} The project delivers an integrated, low-cost architecture for a three-wheeled electric vehicle: fixed-rate dual-sensor acquisition with measured timing and loss, on-board ride features, camera detection fused with time-of-flight ranging, cellular telemetry with local logging and automatic reconnection (replay of the backlog is not yet implemented), and presentation on the vehicle's existing dashboard. Each part has been measured on the bench, and integrating them exposed faults that no subsystem test showed, the most important being the under-voltage caused by powering the LTE modem from the Raspberry Pi. The work has reached a working bench prototype rather than the demonstration in a relevant environment set as the target in Chapter~\ref{chap:intro}; the on-vehicle sessions, and the calibrations listed in Chapter~\ref{chap:discussion}, are what remain to answer RQ1, RQ2 and RQ4 in full. Recommendations for that work are given in Chapter~\ref{chap:recs}.
